\documentclass[10pt,a4paper]{amsart}
\usepackage[margin=19mm]{geometry}
\usepackage{amsmath,amssymb,mathtools}
\usepackage[scr=rsfs,cal=euler]{mathalfa}
\usepackage{xfrac}
\usepackage[unicode,hidelinks,bookmarks=false]{hyperref}
\newcommand{\ie}{i.e.\ }
\newcommand{\cf}{cf.\ }
\newcommand{\resp}{respectively\ }
\newcommand{\Subsection}{\subsection}
\providecommand{\nobreakdash}{\nobreak\hbox{-}}
\setlength{\emergencystretch}{2em}
\multlinegap0pt
\usepackage{amssymb}
\usepackage{mathtools}
\usepackage{tikz}
\usepackage{enumerate}
\usepackage{xcolor}
\newtheorem{corollary}{Corollary}
\newtheorem{lemma}{Lemma}
\title{\bf \Large An improved range for the maximum critically $t$-intersecting hypergraphs}
\author{Lu Lu, Rongrong Lu, Qifan Wang, Tingzeng Wu}
\date{Source: arXiv:2607.28253v1 (CC BY 4.0)}
\begin{document}
\maketitle

\noindent\emph{Source and licence.} \bf \Large An improved range for the maximum critically $t$-intersecting hypergraphs. Version arXiv:2607.28253v1 (CC BY 4.0), distributed by arXiv under the Creative Commons Attribution 4.0 International licence, http://creativecommons.org/licenses/by/4.0/, as recorded in the arXiv licence metadata for this exact version and shown on the abstract page https://arxiv.org/abs/2607.28253v1. Full licence notice: \texttt{licenses/arxiv-2607.28253-CC-BY-4.0.txt}.\par\smallskip

\noindent\textit{Editorial scope.} Counted unit: the article's lemma lem:layer-cake (source0.tex lines 591--611). The article's complete original proof of it is delivered with it (source0.tex lines 612--734, one \texttt{proof} environment of 123 lines). No mathematics is written, repaired or re-derived: the statement and the proof are the article's own lines, in the article's own notation, and no retained line is substituted or re-flowed.  Retained uncounted as the source-local context the proof needs: lines 550--557.  Nothing was omitted from any retained span, so the package carries no omission record.  No retained line contains a citation, so the wrapper carries no bibliography and no bibliography entry.  No retained line contains a figure environment, an image-inclusion command or drawing code, so no image is delivered and the package declares no asset dependency.  Editorial formatting only, in addition: an AMS \texttt{amsart} preamble that restates the article's own declarations and macros used by the retained lines, namely the environments \texttt{corollary}, \texttt{lemma} on separate counters, together with the standard packages those lines need.  The retained environments print in this document as Corollary 1, Lemma 1 in order of appearance, whereas the article prints the same items with the numbering of its own full document.  The complete native source of the article as distributed in the arXiv e-print for this exact version is bundled unchanged as \texttt{sources/206401/source0.tex}.
	\begin{corollary}\label{cor:H-bound}
		Suppose that $k\ge 30d^2$ and $d\ge 3$.  Then, for $0\le s\le \ell-1\le d-1$,
		\[
		|\mathcal H_{\ell,s}|
		\le
		\binom{k-s-1}{\ell-s-1}.
		\]
	\end{corollary}

	\begin{lemma}\label{lem:layer-cake}
		Suppose that $d\ge 3$, $k\ge 30d^2$, and denote
	$
		\lambda:=\frac{8d^2}{k-d}.
	$
		Then we have
		\[
		X
		\le
		\left(
		\frac{2\lambda}{1-\lambda}
		+
		\frac{\lambda^2}{(1-\lambda)^2}
		\right)
		\binom{k}{d}.
		\]
		Consequently,
		\[
		X\le \frac{400}{441}\binom{k}{d}.
		\]
	\end{lemma}

	\begin{proof}
		We first consider the top layer $\ell=d$. 

 Here $a=0$ and $\binom{d}{d}=1$.  Since $M(P)\le d^d\le L_d$, $L_0=1$, a simple stratify gives that
		\[
		(M(P)-1)_+
		\le
		\sum_{s=0}^{d-1}(L_{s+1}-L_s)\mathbf 1_{\{M(P)>L_s\}},
		\]
		where we denote by $\mathbf{1}\{E\}$ the indicator function of an event $E$,
		that is, $\mathbf{1}_{\{E\}} = 1$ if $E$ holds and $0$ otherwise.
		
		Summing over $P\in\binom{E}{d}$ and using Corollary \ref{cor:H-bound}, we get
		\[
		\begin{aligned}
			\sum_{P\in\binom{E}{d}}(M(P)-1)_+
			&\le
			\sum_{s=0}^{d-1}(L_{s+1}-L_s)|\mathcal H_{d,s}|  \\
			&\le
			\sum_{s=0}^{d-1}L_{s+1}\binom{k-s-1}{d-s-1}.
		\end{aligned}
		\]
		Put $r=s+1$.  Since
		$
		\binom{k-r}{d-r}\le \binom{k}{d-r}
		$
		and
		$
		\frac{\binom{k}{d-r}}{\binom{k}{d}}=\prod\limits_{i=0}^{r-1}\frac{d-i}{k-d+i+1}
		\le
		\left(\frac{d}{k-d}\right)^r,
		$
		we have
		\begin{equation}\label{eq:binom-estimate}
		L_r\binom{k-r}{d-r}
		\le
		(8d)^r\left(\frac{d}{k-d}\right)^r\binom{k}{d}
		=
		\lambda^r\binom{k}{d}.
		\end{equation}
		Therefore the contribution of the top layer is at most
		\begin{equation}\label{eq:top-layer}
			\frac{\lambda}{1-\lambda}\binom{k}{d}.
		\end{equation}
		
		Now consider the lower layer $\ell=d-a<d$, where $a\ge 1$.

 Since $M(P)\le d^\ell\le L_\ell$, then
		for each $P\in\binom{E}{\ell}$,
		\[
		M(P)
		\le
		L_a+
		\sum_{\substack{s\ge 0\\ a+s<\ell}}
		(L_{a+s+1}-L_{a+s})\mathbf 1_{\{M(P)>L_{a+s}\}}.
		\]
		Using the simple bound $\left(M(P)-\binom{d}{\ell}\right)_+\le M(P)$, we get the contribution of this layer is at most
		\[
		L_a\binom{k}{d-a}
		+
		\sum_{\substack{s\ge 0\\ a+s<\ell}}
		L_{a+s+1}|\mathcal H_{\ell,s}|.
		\]
	For the second term put
		$
		r=a+s+1.
		$
		By Corollary \ref{cor:H-bound},
		\[
		|\mathcal H_{\ell,s}|
		\le
		\binom{k-s-1}{\ell-s-1}\le \binom{k}{d-r}.
		\]
		Similar to \eqref{eq:binom-estimate},
		$
		L_a\binom{k}{d-a}
		\le
		\lambda^a\binom{k}{d}
		$
		 and $
		L_{a+s+1}|\mathcal H_{\ell,s}|
		\le
		L_r\binom{k}{d-r}
		\le
		\lambda^r\binom{k}{d}.
		$
		Summing over all $a\ge 1$ and over all $r\ge a+1$, we obtain the total lower-layer contribution
		\begin{equation}\label{layer-bound}
		\left(
		\sum_{a\ge 1}\lambda^a
		+
		\sum_{a\ge 1}\sum_{r\ge a+1}\lambda^r
		\right)
		\binom{k}{d}.
		\end{equation}
		Note that $\lambda\le\frac{8d^2}{29d^2}\le\frac{8}{29}<1$, then $\sum\limits_{a\ge 1}\lambda^a=\frac{\lambda}{1-\lambda}$. Besides, $$\sum_{a\ge 1}\sum_{r\ge a+1}\lambda^r=\sum_{a\ge 1}{\frac{\lambda^{a+1}}{1-\lambda}}=\frac{\lambda}{1-\lambda}\sum\limits_{a\ge 1}\lambda^a=\frac{\lambda^2}{(1-\lambda)^2}.$$
		Therefore \eqref{layer-bound} is exactly
		\[
		\left(
		\frac{\lambda}{1-\lambda}
		+
		\frac{\lambda^2}{(1-\lambda)^2}
		\right)
		\binom{k}{d}.
		\]
		Then the total bound is given by combining \eqref{eq:top-layer}.
		
		Finally, since $\lambda\le\frac{8}{29}$,
		then
		$
		\frac{\lambda}{1-\lambda}\le \frac{8}{21}
		$
		and hence
		$
		\frac{2\lambda}{1-\lambda}
		+
		\frac{\lambda^2}{(1-\lambda)^2}
		\le
		\frac{16}{21}+\frac{64}{441}
		=
		\frac{400}{441}.
		$
	\end{proof}

\end{document}
